% appendix/questions/parallelfaster.tex
% mainfile: ../../perfbook.tex
% SPDX-License-Identifier: CC-BY-SA-3.0

\section{为什么并行程序不总是更快？}
\label{sec:app:questions:Why Aren't Parallel Programs Always Faster?}
%
\epigraph{没有什么比简单更复杂的了。}
	 {Charles Poore}

简短的回答是"因为并行执行通常需要通信，而通信不是免费的"。

关于这个问题的更多信息，请参见
\cref{chp:Hardware and its Habits}、
\cref{sec:count:Why Isn't Concurrent Counting Trivial?}，
特别是
\cref{chp:Partitioning and Synchronization Design}，
其中每一章都展示了因拙劣的并行化而使代码变慢的方式。
当然，本书的大部分内容都在讨论如何确保你的并行程序确实比其串行版本更快。

然而，永远不要忘记并行程序可以在相当快的同时也相当简单，
\cref{sec:toolsoftrade:Scripting Languages}
中的示例就是一个很好的例证。
还要记住，并行执行只是众多优化中的一种，
有些程序通过其他优化可能获得更好的效果。
